\section*{अध्याय \ref{ch:b1:stress-strain-rheology} --- प्रतिबल, विकृति और रियोलॉजी}

\begin{solution}{exo:b1:stress-strain-rheology:1}
(क) $\sigma_v = 2700\times9.81\times10\,000 = \qty{2.65e8}{Pa} = \qty{264.9}{MPa}$।
(ख) $P = 1000\times9.81\times10\,000 = \qty{98.1}{MPa}$। (ग) $\sigma_v' =
264.9 - 98.1 = \qty{166.8}{MPa}$, और $\lambda = 98.1/264.9 = 0.370$, जो
$1000/2700$ है।
\end{solution}

\begin{solution}{exo:b1:stress-strain-rheology:2}
(क) $e = (13-10)/10 = 0.30$ और $S = 1.30$। (ख) $e = (7-10)/10 = -0.30$, \qty{30}{\percent} का संकुचन, और $S = 0.70$। (ग) कोण $\psi = \ang{30}$ से सिकुड़ गया है, और $\gamma = \tan\ang{30} = 0.577$।
\end{solution}

\begin{solution}{exo:b1:stress-strain-rheology:3}
(क) केंद्र \qty{70}{MPa}, त्रिज्या \qty{30}{MPa}। (ख) सबसे बड़ा कर्तन-प्रतिबल त्रिज्या है, \qty{30}{MPa}, उन तलों पर जो $\sigma_1$ और $\sigma_3$ से \ang{45} पर हैं, और वे केंद्र का \omterm{def:b1:stress-strain-rheology:stress}{अभिलंब प्रतिबल} वहन करते हैं, \qty{70}{MPa}। (ग)
$2\theta = \ang{60}$: $\sigma_n = 70 + 30\cos\ang{60} = \qty{85}{MPa}$ और
$\tau = 30\sin\ang{60} = \qty{26.0}{MPa}$।
\end{solution}

\begin{solution}{exo:b1:stress-strain-rheology:4}
(क) $\ang{45} - \ang{30}/2 = \ang{30}$ है $\sigma_1$ से, उसके दोनों ओर, अतः दो तल एक-दूसरे से \ang{60} बनाते हैं। (ख) $c' = 20 + r$ के साथ, \cref{thm:b1:stress-strain-rheology:failure-angle} की शर्त पढ़ी जाती है
$r = C\cos\phi + (20 + r)\sin\phi$, अतः
\[
  r = \frac{C\cos\phi + 20\sin\phi}{1 - \sin\phi} = \frac{8.66 + 10}{0.5}
    = \qty{37.3}{MPa},
\]
और $\sigma_1' = 20 + 2r = \qty{94.6}{MPa}$।
\end{solution}

\begin{solution}{exo:b1:stress-strain-rheology:5}
(क) $c = 40$ और $R = \sqrt{20^2 + 15^2} = \qty{25}{MPa}$: $\sigma_1 =
\qty{65}{MPa}$ और $\sigma_3 = \qty{15}{MPa}$; $\tan2\theta_p = 30/40 = 0.75$,
$2\theta_p = \ang{36.87}$, $\theta_p = \ang{18.43}$। (ख) इस कोण पर
$\cos^2\theta = 0.9$, $\sin^2\theta = 0.1$ और $\sin\theta\cos\theta = 0.3$:
$\sigma_n = 60\times0.9 + 2\times15\times0.3 + 20\times0.1 = 54 + 9 + 2 =
\qty{65}{MPa}$, और $\tau = (20-60)\times0.3 + 15\times(0.9-0.1) = -12 + 12 =
0$। (ग) $R = \qty{25}{MPa}$, मुख्य दिशाओं से \ang{45} पर के तलों पर।
\end{solution}

\begin{solution}{exo:b1:stress-strain-rheology:6}
(क) $S_{1,2} = \sqrt{1 + 1} \pm 1$: $S_1 = 2.414$ और $S_2 = 0.414$;
$\tan2\theta' = 2/2 = 1$, $\theta' = \ang{22.5}$। (ख) $S_1S_2 =
(\sqrt2 + 1)(\sqrt2 - 1) = 1$, जो $\det F$ है: कोई क्षेत्रफल-परिवर्तन नहीं। (ग) $(1, 1)$ के अनुदिश रेखा $F(1, 1) = (1 + 2, 1) = (3, 1)$ पर जाती है, कर्तन-दिशा से $\arctan(1/3) = \ang{18.4}$ पर, \omterm{def:b1:stress-strain-rheology:strain}{स्ट्रेच}
$\sqrt{10}/\sqrt2 = 2.24$ के साथ। वह दीर्घवृत्त के लंबे अक्ष से, जो \ang{22.5} पर है, आगे निकल गई है, कर्तन-तल की ओर; उसका \omterm{def:b1:stress-strain-rheology:strain}{स्ट्रेच}
$S_1$ से कुछ कम है, जैसा होना ही चाहिए, क्योंकि कोई रेखा $S_1$ से अधिक नहीं खिंचती।
\end{solution}

\begin{solution}{exo:b1:stress-strain-rheology:7}
(क) $S_1/S_2 = 4$ और $S_1S_2 = 1$ से $S_1 = 2$ और $S_2 = 0.5$ मिलते हैं;
\cref{prop:b1:stress-strain-rheology:simple-shear} द्वारा, $\gamma = S_1 - S_2 =
1.5$। (ख) $\tan2\theta' = 2/1.5 = 1.333$, $2\theta' = \ang{53.1}$: लंबे
अक्ष दीवारों से \ang{26.6} पर होने चाहिए। (ग) $1.5\times200 = \qty{300}{m}$।
\end{solution}

\begin{solution}{exo:b1:stress-strain-rheology:8}
$\sigma_v = 2700\times9.81\times3000 = \qty{79.46}{MPa}$। (क) $P = \qty{29.43}{MPa}$,
$\sigma_n' = 70 - 29.43 = \qty{40.57}{MPa}$, और बायरली अनुमति देता है
$0.85\times40.57 = \qty{34.5}{MPa}$, जो लागू \qty{30}{MPa} से अधिक है: भ्रंश टिकता है। (ख) $P = 0.6\times79.46 = \qty{47.68}{MPa}$, $\sigma_n' =
\qty{22.32}{MPa}$, और घर्षण केवल \qty{19.0}{MPa} देता है: भ्रंश
फिसलता है। (ग) $30 = 0.85\,(70 - 79.46\,\lambda)$ से $\lambda = (70 -
35.29)/79.46 = 0.437$।
\end{solution}

\begin{solution}{exo:b1:stress-strain-rheology:9}
(क) प्रतिबल $\dot\varepsilon^{1/4}$ के अनुसार बढ़ता है: $10^{1/4} = 1.78$। (ख) सामर्थ्यों का
अनुपात है
\[
  \exp\Bigl[\frac{E}{nR}\Bigl(\frac{1}{673.15} - \frac{1}{773.15}\Bigr)\Bigr]
  = \exp\bigl(6705\times\num{1.921e-4}\bigr) = \exp(1.288) = 3.63 .
\]
(ग) दर का दसगुना परिवर्तन विसर्पण-सामर्थ्य को 1.78 गुणक से बदलता है,
जो तापमान केवल $100\times\ln1.78/\ln3.63 \approx
\qty{45}{\degreeCelsius}$ में करता है; लगभग \qty{20}{\degreeCelsius} प्रति किलोमीटर के \omterm{def:b1:heat-budget:geotherm}{जियोथर्म} पर, वह संक्रमण को लगभग \qty{2}{km} खिसकाता है। \omterm{def:b1:heat-budget:geotherm}{जियोथर्म} का एक परिवर्तन समतापों को, और उनके साथ संक्रमण को, इससे कहीं अधिक खिसकाता है।
\end{solution}

\begin{solution}{exo:b1:stress-strain-rheology:10}
(क) $C = 0$ के साथ, भ्रंश तब फिसलता है जब $r = \mu c'/\sqrt{1+\mu^2}$
(\cref{thm:b1:stress-strain-rheology:failure-angle})। यहाँ $c' = S =
\rho g z(1-\lambda)$ है, अतः $\sigma_1 - \sigma_3 = 2r = 2\mu S/\sqrt{1+\mu^2}$।
(ख) $S$ प्रति किलोमीटर $2700\times9.81\times1000\times(1 - 1000/2700) =
\qty{16.68}{MPa}$ बढ़ता है, और $2\times0.85/1.3124 = 1.295$: ढलान
प्रति किलोमीटर \qty{21.6}{MPa} है, \num{13.1} (सामान्य भ्रंश) और
\qty{61.3}{MPa} प्रति किलोमीटर (उत्क्षेप) के बीच। (ग) $2\mu/(k+\mu) = 0.786$,
$2\mu/k = 1.295$ और $2\mu/(k-\mu) = 3.68$, $k = 1.3124$ के साथ: फिसलने पर \omterm{def:b1:stress-strain-rheology:principal}{अवकल प्रतिबल}
ऊर्ध्वाधर \omterm{def:b1:stress-strain-rheology:lithostatic}{प्रभावी प्रतिबल} का 0.79, 1.30 और 3.68 गुना है।
\end{solution}

\begin{solution}{exo:b1:stress-strain-rheology:11}
(क) $t = 60\times\num{e6}\times\num{3.156e7} = \qty{1.89e15}{s}$ और
$2\sqrt{\kappa t} = 2\sqrt{\num{1.89e9}} = \qty{87.0}{km}$। (ख)
$\operatorname{erf}x = 600/1350 = 0.444$ के लिए $x = 0.417$: $z = 0.417\times87.0
= \qty{36.3}{km}$; $\operatorname{erf}x = 800/1350 = 0.593$ के लिए $x = 0.586$:
$z = \qty{51.0}{km}$। (ग) भूकंप \qty{36}{km} से \qty{51}{km} तक पहुँचने चाहिए;
आवरण भंगुर परत का आधार \qty{38.2}{km} पर रखता है, उस परास के भीतर और उसके ऊपरी भाग के निकट, \qty{641}{\degreeCelsius} के तापमान पर, जो देखे गए परास में है।
\end{solution}

\begin{solution}{exo:b1:stress-strain-rheology:12}
(क) $r = 45$ और $c' = \qty{65}{MPa}$; केंद्र से रेखा $\tau = 0.85\sigma_n'$ तक की दूरी $0.85\times65/1.3124 = \qty{42.1}{MPa}$ है, त्रिज्या से कम:
वृत्त रेखा को पार करता है, और उसके ऊपर के चाप पर के तल फिसलते हैं। (ख) $2\theta$ को $\sigma_1$ सिरे से कोण मानते हुए, फिसलने के लिए $45\sin2\theta \ge 0.85\,(65 + 45\cos2\theta)$ चाहिए, अर्थात्
$45\sin2\theta - 38.25\cos2\theta \ge 55.25$, या
$59.06\sin(2\theta - \ang{40.36}) \ge 55.25$, जहाँ $\ang{40.36} = \phi$।
अतः $\sin(2\theta - \phi) \ge 0.9355$, $2\theta - \phi$ \ang{69.31} और
\ang{110.69} के बीच, $2\theta$ \ang{109.67} और \ang{151.06} के बीच: अभिलंब
$\sigma_1$ से \ang{54.8} से \ang{75.5} पर हैं, और फिसलने वाले भंग
$\sigma_1$ से \ang{14.5} से \ang{35.2} पर। परास का मध्य,
\ang{24.8}, $\ang{45} - \phi/2$ है, वह तल जो
\cref{thm:b1:stress-strain-rheology:failure-angle} का है। (ग) केवल वही तल फिसलता है जब वृत्त रेखा को बस छूता है:
$r = \mu(20 + r)/k$, $r = 20\mu/(k -
\mu) = 17/0.4624 = \qty{36.8}{MPa}$, और $\sigma_1' = 20 + 2r =
\qty{93.5}{MPa}$।
\end{solution}

\begin{solution}{pb:b1:stress-strain-rheology:1}
\textbf{1.} $\sigma_v = 2700\times9.81\times5000 = \qty{132.44}{MPa}$।

\textbf{2.} $P = 1000\times9.81\times5000 = \qty{49.05}{MPa}$, और
$\lambda = 49.05/132.44 = 0.370$।

\textbf{3.} $\sigma_1 = \sigma_v = \qty{132.44}{MPa}$, ऊर्ध्वाधर, और
$\sigma_3 = \qty{75.00}{MPa}$, क्षैतिज पूर्व–पश्चिम: \omterm{def:b1:layered-earth:layers}{भूपर्पटी} अपने ही भार के नीचे पूर्व–पश्चिम खिंच रही है, वह अवस्था जो सामान्य भ्रंश चलाती है, जैसा यह है।

\textbf{4.} $\sigma_1 - \sigma_3 = \qty{57.44}{MPa}$; केंद्र
$(132.44 + 75.00)/2 = \qty{103.72}{MPa}$; त्रिज्या \qty{28.72}{MPa}।

\textbf{5.} $\sigma_1' = 132.44 - 49.05 = \qty{83.39}{MPa}$, $\sigma_3' =
75.00 - 49.05 = \qty{25.95}{MPa}$, केंद्र $103.72 - 49.05 = \qty{54.67}{MPa}$।

\textbf{6.} $\sigma_{xx} = \qty{75.00}{MPa}$, $\sigma_{yy} = \qty{132.44}{MPa}$
और $\sigma_{xy} = 0$: अक्ष मुख्य दिशाएँ हैं, उनमें से एक ऊर्ध्वाधर, अतः उन पर लंबवत तल कोई कर्तन-प्रतिबल नहीं वहन करते।

\textbf{7.} भ्रंश अपनी \omterm{def:b1:geological-maps:attitude}{नति} के अनुदिश इकाई सदिश
$(\cos\ang{60}, -\sin\ang{60})$ के अनुदिश चलता है, पूर्व और नीचे की ओर; उसका ऊपर का अभिलंब
$\vec n = (\sin\ang{60}, \cos\ang{60}) = (0.866, 0.500)$ है, $x$ से \ang{30} पर। $\sigma_1$ के साथ उसके कोण की, $y$ के अनुदिश, कोज्या
0.500 है: $\theta = \ang{60}$।

\textbf{8.} $x$ से \ang{30} पर: $\sigma_n = 75.00\times0.75 + 132.44\times0.25
= 56.25 + 33.11 = \qty{89.36}{MPa}$, और $\tau = (132.44 - 75.00)\times
\sin\ang{30}\cos\ang{30} = 57.44\times0.4330 = \qty{24.87}{MPa}$।

\textbf{9.} $2\theta = \ang{120}$ पर, $\sigma_1$ सिरे से:
$\sigma_n = 103.72 + 28.72\cos\ang{120} = 103.72 - 14.36 = \qty{89.36}{MPa}$
और $\tau = 28.72\sin\ang{120} = \qty{24.87}{MPa}$।

\textbf{10.} $\sigma_n' = 89.36 - 49.05 = \qty{40.31}{MPa}$।

\textbf{11.} $24.87/40.31 = 0.617$, 0.85 से नीचे: भ्रंश टिकता है।

\textbf{12.} $0.85\times40.31 = \qty{34.26}{MPa}$; कर्तन-प्रतिबल उससे
\qty{9.39}{MPa} कम है।

\textbf{13.} वृत्त का केंद्र \qty{54.67}{MPa} पर है और त्रिज्या
\qty{28.72}{MPa}, \num{25.95} से \qty{83.39}{MPa} तक; बायरली रेखा
$\tau = 0.85\,\sigma_n'$ है (यहाँ सभी \omterm{def:b1:stress-strain-rheology:stress}{अभिलंब प्रतिबल} \qty{200}{MPa} से नीचे हैं); भ्रंश बिंदु $(40.31, 24.87)$ है, $\sigma_1$ सिरे से \ang{120} पर।

\textbf{14.} $0.85\times54.67/\sqrt{1 + 0.85^2} = 46.47/1.3124 =
\qty{35.41}{MPa}$, त्रिज्या \qty{28.72}{MPa} से अधिक: पूरा वृत्त रेखा के नीचे है, और चट्टान से होकर कोई तल फिसल नहीं सकता।

\textbf{15.} $\phi = \arctan0.85 = \ang{40.36}$: तल $\ang{45} - \phi/2 = \ang{24.8}$ पर है, ऊर्ध्वाधर $\sigma_1$ से, \ang{65.2} की \omterm{def:b1:geological-maps:attitude}{नति} से।

\textbf{16.} $S = \sigma_v' = \qty{83.39}{MPa}$ और
$\sigma_1 - \sigma_3 = 2\times0.85\times83.39/(1.3124 + 0.85) =
\qty{65.55}{MPa}$; ढलान $13.11\times5 = \qty{65.55}{MPa}$ देती है।

\textbf{17.} $\sigma_3 = 132.44 - 65.55 = \qty{66.88}{MPa}$।

\textbf{18.} भ्रंश तब फिसलता है जब
$r\sin\ang{120} = 0.85\,(c' + r\cos\ang{120})$ के साथ
$c' = (83.39 + \sigma_3')/2$ और $r = (83.39 - \sigma_3')/2$; हल करने पर
$\sigma_3' = \qty{17.18}{MPa}$, $\sigma_3 = \qty{66.23}{MPa}$। यह सर्वोत्तम तल से कम है: भ्रंश उस बिंदु पर नहीं है जहाँ वृत्त पहले रेखा को छूता है, अतः वृत्त को भ्रंश के अपने बिंदु तक पहुँचने से पहले और बढ़ना होगा।

\textbf{19.} प्रत्येक प्रभावी \omterm{def:b1:stress-strain-rheology:stress}{अभिलंब प्रतिबल}, $\sigma_1'$ और $\sigma_3'$ सहित, $\Delta P$ से घटता है, और उनका अंतर नहीं बदलता; कर्तन-प्रतिबल छिद्र-दाब पर बिल्कुल निर्भर नहीं करते।

\textbf{20.} सर्वोत्तम तल तब फिसलता है जब दूरी त्रिज्या के बराबर हो:
$0.85\,(54.67 - \Delta P)/1.3124 = 28.72$, अतः
$\Delta P = 54.67 - 28.72\times1.3124/0.85 = 54.67 - 44.34 = \qty{10.33}{MPa}$।

\textbf{21.} कूप में पानी का स्तंभ पहले जितना भारी है, अतः पंप को कूप-मुख पर पूरा \qty{10.33}{MPa} जोड़ना होगा, वायुमंडलीय दाब का लगभग 103 गुना:
पानी के स्तंभ का भार $\num{10.33e6}/(1000\times9.81) = \qty{1053}{m}$ ऊँचा।

\textbf{22.} $\mu = 0.6$ के साथ भ्रंश केवल
$0.6\times40.31 = \qty{24.19}{MPa}$ वहन कर सकता, जो उससे कम है जो वह वहन करता है, \qty{24.87}{MPa}: वह पहले से फिसल रहा होता। चूँकि वह शांत है, उसका घर्षण कम से कम $0.617$ है, या उसमें कुछ \omterm{def:b1:stress-strain-rheology:coulomb}{संसंजन} है।

\textbf{23.} सर्वोत्तम तल तब फिसलता है जब छिद्र-दाब पहुँचे $49.05 + 10.33 = \qty{59.38}{MPa}$, अर्थात् $\lambda = 59.38/132.44 = 0.448$, स्थैतिक मान 1 से काफ़ी नीचे।

\textbf{24.} $\Delta P = \sigma_n' - \tau/0.85 = 40.309 - 24.870/0.85 = 40.309 -
29.259 = \qty{11.049}{MPa}$, सर्वोत्तम \omterm{def:b1:geological-maps:attitude}{अभिविन्यास} के तल को हिलाने वाली \qty{10.33}{MPa} से अधिक। एक दशमलव तक, भ्रंश को फिसलाने वाली छिद्र-दाब-वृद्धि है
\textbf{\qty[text-series-to-math=true]{11.0}{MPa}}।
\end{solution}
